\section{The model and the monotone envelope}
\label{sec:model}

There are $N$ fixed owners, indexed by $I$, and $D$ production resources.
A report is a vector $b_i\in\mathbb R_+^D$.  The production function is
\begin{equation}
 F(y)=\min_{\ell\in\mathcal L} w^\ell\cdot y,
 \qquad w^\ell\in\mathbb R_+^D,
 \label{eq:minlinear}
\end{equation}
where $\mathcal L$ is finite and nonempty.  Thus $F$ is nonnegative,
monotone, concave, and positively homogeneous on the entire nonnegative
orthant.  At a profile $\mathbf b=(b_i)_{i\in I}$, write
\[
 v_{\mathbf b}(S)=F\Bigl(\sum_{i\in S}b_i\Bigr),
 \qquad T(\mathbf b)=v_{\mathbf b}(I).
\]
The coalition-cover polyhedron and efficient core are, respectively,
\[
 \mathcal P(v)=\{x\in\mathbb R_+^I:x(S)\ge v(S)\text{ for all }S\subseteq I\},
 \qquad
 \mathcal C(v)=\{x\in\mathcal P(v):x(I)=v(I)\}.
\]
Here $x(S)=\sum_{i\in S}x_i$.  The core is nonempty: choose a form
$w^\ell$ attaining the minimum in \eqref{eq:minlinear} at the aggregate
report and set $x_i=w^\ell\cdot b_i$.  This vector covers every coalition
and has total $T(\mathbf b)$.  Conversely, every cover has total at least
$T(\mathbf b)$ by the grand-coalition inequality.

A payment rule is \emph{own-monotone} if
\[
 M_i(a_i,b_{-i})\le M_i(b_i,b_{-i})\qquad(0\le a_i\le b_i),
\]
where vector inequalities are coordinatewise.  This is the individual
withholding condition studied here.  In particular, no restriction to
Boolean reports is part of the model.  An owner's utility is its received
payment: there is no additional outside value for withheld resources or
private cost term.  All funding ratios below are total payment divided by
output; the corresponding net external subsidy ratio is one less.

\begin{definition}[Seed and true envelope]
A seed $q$ selects a vector in $\mathcal P(v_{\mathbf b})$ at every real
report profile.  Its true own-report envelope is
\begin{equation}
 E_i[q](\mathbf b)=\sup_{0\le a_i\le b_i}q_i(a_i,b_{-i}).
 \label{eq:envelope}
\end{equation}
The supremum initially takes values in $[0,\infty]$.
\end{definition}

\begin{lemma}[Least monotone majorant]
\label{lem:least-majorant}
The envelope $E[q]$ is own-monotone, dominates $q$, and covers every
coalition.  If an own-monotone payment rule $M$ satisfies
$M_i(\mathbf a)\ge q_i(\mathbf a)$ for every owner and every profile, then
$M\ge E[q]$ coordinatewise.  Consequently $E[q]$ is the coordinatewise
least own-monotone majorant of $q$.
\end{lemma}
\begin{proof}
The set over which the supremum in \eqref{eq:envelope} is taken expands
when $b_i$ increases, proving own-monotonicity.  The choice $a_i=b_i$
gives $E\ge q$, which also proves coalition coverage.  For every admissible
$a_i$, own-monotonicity and the majorant condition give
$M_i(\mathbf b)\ge M_i(a_i,b_{-i})\ge q_i(a_i,b_{-i})$.
Taking the supremum proves the last assertion.
\end{proof}

The majorant condition in Lemma~\ref{lem:least-majorant} is an additional
baseline-payment floor, imposed at \emph{every} profile.  Our lower bounds
therefore apply to all monotone repairs that only add to the specified
seed payments.  They do not impose such a floor on unrestricted mechanisms,
and they do not establish such growing bounds for every own-monotone cover rule.
The below-two theorem does have the separate constant-factor consequence
stated in Remark~\ref{rem:constant-two}.

For $C\ge1$ and $r>1$, define the capped power seed
\begin{equation}
 q^{C,r}(\mathbf b)=
 \mathop{\rm argmin}_{x\in\mathcal P(v_{\mathbf b}),\ x(I)\le C T(\mathbf b)}
       \sum_{i\in I}x_i^r.
 \label{eq:capped-seed}
\end{equation}
The feasible set is nonempty and compact, and strict convexity makes the
minimizer unique.  When $T=0$ it is zero.  We also use the uncapped quadratic
seed
\begin{equation}
 q^{\infty,2}(\mathbf b)=\mathop{\rm argmin}_{x\in\mathcal P(v_{\mathbf b})}
       \|x\|_2^2.
 \label{eq:uncapped-seed}
\end{equation}
It exists and is unique.  An efficient cover $c$ has
$\|c\|_2\le c(I)=T$, so its projection satisfies
$q_i^{\infty,2}\le\|q^{\infty,2}\|_2\le T$.

\begin{proposition}[Elementary funding upper bounds]
\label{prop:trivial-upper}
If a seed satisfies $q(I;\mathbf b)\le C T(\mathbf b)$ at every profile,
then
\[
 E_i[q](\mathbf b)\le C T(\mathbf b),\qquad
 E[q](I;\mathbf b)\le NC T(\mathbf b).
\]
For the uncapped quadratic seed the same statements hold with $C=1$.
\end{proposition}
\begin{proof}
For every underreport, monotonicity of $F$ gives
$T(a_i,b_{-i})\le T(\mathbf b)$.  Apply the assumed budget bound, or the
coordinate bound for \eqref{eq:uncapped-seed}, before taking the supremum.
\end{proof}

\section{Efficient seeds and seeds with budget below two}
\label{sec:pool-barriers}

An elementary pool technology already gives a linear envelope obstruction
for every efficient-core selector.  This statement does not require a
particular objective, symmetry, or continuity of the selector.

\begin{theorem}[Every efficient-core selector]
\label{thm:core-barrier}
For every integer $g\ge2$, there is one production function of the form
\eqref{eq:minlinear}, with binary coefficients, $N=D=2g$ owners and
resources, and $g$ designated profiles, such that every seed selecting an
efficient core vector at these profiles and at their designated unilateral
underreports has
\begin{equation}
 \frac{E[q](I;\mathbf b)}{T(\mathbf b)}
 \ge \frac{g+3}{4}=\frac{N+6}{8}
 \label{eq:core-bound}
\end{equation}
at at least one designated profile.  The bad profile may depend on the
selector.
\end{theorem}
\begin{proof}
Partition the owners into $g$ pairs $(h_j,l_j)$, each with its own scalar
resource, and put
$F(y)=\min_j(y_{h_j}+y_{l_j})$.
At profile $F_j$, pool $j$ has capacities $(1,1)$ and every other pool has
capacities $(4,1)$.  At profile $L_{ij}$, pools $i,j$ have capacities
$(1,1)$ and the remaining pools have capacities $(4,1)$.  Every such
profile has output $2$.  The profile $L_{ij}$ is obtained from $F_j$ solely
by owner $h_i$ reducing its capacity from $4$ to $1$.

At $F_j$, the coalition containing both small-pool owners and every
large-pool high owner has value $2$.  Efficiency therefore forces all
large-pool low owners to receive zero.  A coalition consisting of either
small-pool owner and all large-pool low owners has value $1$, forcing each
small-pool owner to receive at least $1$.  Thus the core is the singleton
that pays these two owners $1$ each.

At $L_{ij}$, the same argument first makes all remaining large-pool low
payments zero.  Each coalition containing one owner from each small pool
and all those low owners has value $1$.  Two disjoint choices force the
four small-pool payments to total at least $2$, so all other payments
vanish.  The four cross-pair inequalities and efficiency then give exactly
\[
 (x_{h_i},x_{l_i};x_{h_j},x_{l_j})
       =(\lambda,\lambda;1-\lambda,1-\lambda),\quad 0\le\lambda\le1.
\]
Conversely every displayed vector covers: its payment to a coalition is
a convex combination of that coalition's owner counts in the two small
pools, and hence is at least their minimum, which bounds production.

Set $a_{ij}=q_{h_i}(L_{ij})$.  Then $a_{ij}+a_{ji}=1$, and at $F_j$
the envelope total is at least $2+\sum_{i\ne j}a_{ij}$.  Averaging over
$j$ gives a lower bound $2+(g-1)/2$ at some $F_j$.  Divide by its output
$2$ to obtain \eqref{eq:core-bound}.
\end{proof}

\begin{remark}
The games used here need not be convex cooperative games.  At $F_j$, let
$S$ contain $h_j$ and every other high owner, and let $U$ contain $h_j$
and every other low owner.  Then $v(S)=v(U)=v(S\cup U)=1$ and
$v(S\cap U)=0$, violating supermodularity.  The construction therefore
does not conflict with positive selection results restricted to convex
games.
\end{remark}

\begin{theorem}[Every seed with budget factor below two]
\label{thm:seed-barrier}
Fix $1\le B<2$.  Choose an integer $m\ge2$ with $m>1/(2-B)$ and set
\[
 \gamma=1-\frac{(B-1)m}{m-1}>0.
\]
For each $g\ge2$ there is a binary finite-minlinear technology with
$N=D=gm$ and finitely many profiles of common output $m$ such that every
cover seed whose total is at most $Bm$ on these profiles satisfies, at
one of the full designated profiles,
\begin{equation}
 \frac{E[q](I;\mathbf b)}{T(\mathbf b)}
 \ge \gamma\left(1+\frac{g-1}{2m}\right).
 \label{eq:seed-below-two}
\end{equation}
In particular, every globally $B$-funded seed has a linear envelope
obstruction along this sequence of owner counts.
\end{theorem}
\begin{proof}
Use $g$ pools of $m$ scalar owners and let $F$ be the minimum pool total.
Each pool has a designated high owner and $m-1$ low owners.  A small pool
has all capacities $1$; a large pool has high capacity $R\ge m$ and all
low capacities $1$.  Define $F_j$ to have just pool $j$ small, and
$L_{ij}$ to have just pools $i,j$ small.  These profiles have output $m$,
and $L_{ij}$ is an own underreport by $h_i$ from $F_j$.

At either type of profile, let $L_*$ denote the sum of seed payments to
the low owners in all large pools.  The coalition of all small-pool
owners and all large-pool high owners has value $m$.  Hence the seed's
budget implies $L_*\le(B-1)m$.  At $F_j$, any one small-pool owner together
with one chosen low owner from each large pool has value $1$.  Averaging
over the $m-1$ equally likely low choices in each large pool gives
$q_a+L_* /(m-1)\ge1$.  Thus the small-pool payments total at least
$m\gamma$.

At $L_{ij}$, the coalition of $h_i,h_j$ and one low owner from each
remaining large pool has value $1$.  The same averaging gives
$q_{h_i}(L_{ij})+q_{h_j}(L_{ij})\ge\gamma$.
Consequently the envelope at $F_j$ is at least
$m\gamma+\sum_{i\ne j}q_{h_i}(L_{ij})$.  Summing this inequality over
$j$ and then averaging proves \eqref{eq:seed-below-two}.  Empty collections
of large pools cause no change to these inequalities.
\end{proof}

\begin{remark}[A constant-factor consequence without a baseline floor]
\label{rem:constant-two}
For each fixed $1\le B<2$, choose $m$ as in Theorem~\ref{thm:seed-barrier}
and then $g$ so large that $\gamma[1+(g-1)/(2m)]>B$.  On the resulting one
fixed technology, no own-monotone covering rule $M$ can satisfy
$M(I;\mathbf b)\le B T(\mathbf b)$ at every profile.  Indeed,
$E[M]=M$ by own-monotonicity, so taking $q=M$ contradicts the theorem.
This excludes a universal funding factor below two for unrestricted rules.
It does not give a growing owner-count lower bound for unrestricted rules
whose funding factors may exceed two.  We record it as an immediate
consequence, not an independent contribution.
\end{remark}

\begin{corollary}[The linear-objective endpoint]
\label{cor:r-one}
For any cap $C\ge1$, every tie selection from the minimizers of
$\sum_i x_i$ over $\mathcal P(v_{\mathbf b})$ with total at most $CT$
obeys Theorem~\ref{thm:core-barrier}.  The same is true at cap $C=1$
for every objective and every feasible selector.
\end{corollary}
\begin{proof}
The minimum possible total is exactly $T$, and its minimizers are precisely
the efficient core.  A cap of $T$ likewise restricts all feasible vectors
to that core.
\end{proof}

\section{Every fixed capped power seed has linear inflation}
\label{sec:fixed-caps}

The next result includes the quadratic seed with cap $2T$.  The proof uses
a finite incidence construction; its complete distributional and production
realization arguments are given in Appendices~\ref{app:qary} and
\ref{app:realization}.  The integer $L$ below is a local combinatorial
dimension, whereas $D$ always denotes the number of production resources.

\begin{theorem}[Fixed caps and fixed powers]
\label{thm:fixed-cap}
Fix $C>1$ and $r>1$.  There is a constant $c_C>0$ and a sequence of
binary finite-minlinear production instances with unbounded owner count
$N$ such that
\[
 E[q^{C,r}](I;\mathbf b)\ge c_C(N-2)T(\mathbf b).
\]
The constant can be chosen independently of $r$, although the first
admissible instance may depend on $r$.  More explicitly, for any integer
$Q\ge3$ with
\begin{equation}
 C_Q=\frac{(Q-1)L+2h}{L+h}
     =\frac{2Q^2+2Q-2}{2Q+5}>C,
 \qquad \frac hL=\frac1{2(Q+2)},
 \label{eq:natural-cap}
\end{equation}
one may take $c_C=(C-1)/(2Q^2-7)$, with $D=2N-1$.
For $C=2$, no added owner is needed: along a sequence with $D=2N$,
\begin{equation}
 E[q^{2,r}](I;\mathbf b)\ge\frac{N-1}{11}T(\mathbf b).
 \label{eq:cap-two-eleventh}
\end{equation}
For $r=2$ every local dimension in the construction is admissible.
At the remaining endpoint $C=1$, Corollary~\ref{cor:r-one} and
Theorem~\ref{thm:core-barrier} give the bound $(N+6)/8$ for every $r>1$
along the separate sequence $N=D=2g$; the formula for $c_C$ above is
not used at that endpoint.
Together with Proposition~\ref{prop:trivial-upper}, these are linear-order
upper and lower bounds for each fixed capped power architecture.
\end{theorem}

We first state the combinatorial input with all of its parameters.

\begin{lemma}[A transitive row family with focal mixtures]
\label{lem:qary-input}
Fix integers $Q\ge3$, $j\ge1$, and put
\begin{align*}
 L&=16jQ(Q-1)(Q+2),& h&=8jQ(Q-1),& \xi&=h/L,\\
 n&=Q^L,& m&=n-(Q-1)L-1,&
 \mu&=\frac{Q-2+\xi}{Q-1},\\
 \delta&=\frac1{4Q(Q-1)},& \beta_0&=\mu+\delta,&
 K&=\lfloor\beta_0m\rfloor,\quad \beta=K/m,\\
 s&=(Q-2)L+h+K,& k&=n-s,& a&=1-\mu,\quad b_*=1-\beta.
\end{align*}
On $V=(\mathbb Z/Q\mathbb Z)^L$, let $\mathcal R$ consist of the
complements of all independent $s$-sets for the hypergraph of complete
axis-parallel $Q$-point lines.  Here independent means containing no
complete line.  This is a nonempty translation-invariant family of
$k$-element rows, and
\begin{equation}
 \frac nQ\le k<\frac n{Q-1},\quad
 a-b_*>\frac\delta2,\quad
 h(a-b_*)>1-b_*,\quad b_*>\frac1Q.
 \label{eq:qary-margins}
\end{equation}
For each focal vertex $i$ there is a finite rational probability mixture
of rows containing $i$ whose mean $w^i$ has coordinate $1$ at $i$,
coordinate $a$ at each of its $(Q-1)L$ axis neighbors, and coordinate
$b_*$ elsewhere.  Every supported row contains exactly $L-h$ neighbors
and $m-K$ other vertices.  Every row not containing $i$ contains at least
$L$ neighbors of $i$.
\end{lemma}

For a focal vertex $i$, retain demand $1$ on every off-$i$ row and on
the support of its mixture; replace demand by $0$ on all other rows.
Adjoin one common veto owner $u$ to each row.  The resulting full and
focal games are precisely realized by one all-real production function
with $n+1$ owners and $2(n+1)$ resources, by
Lemma~\ref{lem:realization}.  All designated profiles have output $1$.

\begin{lemma}[The natural-cap certificate]
\label{lem:power-certificate}
Use Lemma~\ref{lem:qary-input}.  Suppose $r>1$ and
\begin{equation}
 h^{r-1}(a-b_*)\ge1-b_*.
 \label{eq:power-condition}
\end{equation}
At the focal underreport, the unique minimizer for cap $C_Q$ pays
\[
 q_u=q_i=p=\frac h{L+h},\qquad
 q_v=z=\frac1{L+h}\quad(v\text{ an axis neighbor of }i),
 \qquad q_v=0\quad\text{otherwise}.
\]
For each fixed $r>1$, condition \eqref{eq:power-condition} holds for all
sufficiently large $j$; for $r=2$ it holds for every $j\ge1$.
\end{lemma}
\begin{proof}
Every retained own row has payment $2p+(L-h)z=1$, and every off row has
payment at least $p+Lz=1$.  Deleted rows have demand zero.  The total
payment is exactly $C_Q$.

Use the equivalent objective $\Phi(x)=\sum_v x_v^r/r$.  Give the supported
own rows total multiplier
\[
 \lambda=\frac{p^{r-1}-z^{r-1}}{1-a}>0
\]
distributed according to their probability mixture, and set
$\eta=\lambda-p^{r-1}$.  A direct calculation shows
\[
 \eta-\lambda b_*\ge0
 \quad\Longleftrightarrow\quad
 h^{r-1}(a-b_*)\ge1-b_*.
\]
In particular $\eta\ge0$.  The gradient of $\Phi$ at the proposed vector
is $\lambda$ times the mean augmented row minus $\eta\mathbf1$, with
the nonnegative correction $\eta-\lambda b_*$ on the zero coordinates.
For any feasible $y$, nonnegativity, the retained row constraints, and
$y(I)\le C_Q$ therefore give
\[
 \nabla\Phi(q)\cdot y\ge\lambda-\eta C_Q
                          =\nabla\Phi(q)\cdot q.
\]
Convexity proves optimality, and strict convexity proves uniqueness.
This argument does not require a constraint qualification.

By \eqref{eq:qary-margins}, $a-b_*>\delta/2$, so
$h^{r-1}\ge2/\delta$ suffices.  This holds eventually for fixed $r>1$.
For $r=2$, the sharper inequality $h(a-b_*)>1-b_*$ already holds for
every parameter value.
\end{proof}

\begin{proof}[Proof of Theorem~\ref{thm:fixed-cap}]
For $Q=3$, \eqref{eq:natural-cap} gives $C_Q=2$ and $p=1/11$.
Translate the focal certificate to every one of the $n$ private owners
and apply the simultaneous realization.  Each owner has an admissible
own underreport paying $1/11$, proving \eqref{eq:cap-two-eleventh}.

For general $C>1$, choose $Q$ with $C_Q>C$; such $Q$ exist since
$C_Q$ is unbounded.  Append one resource and one owner with fixed amount
\[
 A=\frac{C_Q-C}{C-1}>0,
\]
and replace the old technology by $F_{\rm old}+Y$.  Its min-linear forms
still have binary coefficients.  At the full profile and every selected
old-owner underreport, total output is $1+A$.  The cover constraints
factor into the old cover constraints and $x_{\rm new}\ge A$; these
conditions also suffice for mixed coalitions.  Since the objective is
strictly increasing in this coordinate above $A$, its minimizer has
$x_{\rm new}=A$.  The residual old-owner cap is
$C(1+A)-A=C_Q$, so Lemma~\ref{lem:power-certificate} applies unchanged.
There are $N=n+2$ owners, and
\[
 \frac{np}{1+A}
 =\frac{C-1}{2Q^2-7}(N-2).
\]
The resource count is $2(n+1)+1=2N-1$.  All real reports remain in the
domain; only the selected reports are used for the lower bound.
\end{proof}

\begin{corollary}[Repairs preserving the seed]
\label{cor:majorant-barrier}
Every own-monotone rule that dominates one of the seeds in
Theorems~\ref{thm:core-barrier}, \ref{thm:seed-barrier}, or
\ref{thm:fixed-cap} at every profile inherits that theorem's funding lower
bound.  For the specified capped seeds this applies to any repair by
top-ups, independently of how the repair is implemented.
\end{corollary}
\begin{proof}
Apply Lemma~\ref{lem:least-majorant} at the designated bad profile.
\end{proof}

The parameters $C$ and $r$ in Theorem~\ref{thm:fixed-cap} are fixed before
the instance grows.  The theorem does not assert a starting dimension
uniform in $r$, a fixed production-resource dimension, or an obstruction
to an unrestricted own-monotone mechanism.

\section{The uncapped quadratic seed exceeds every square-root constant}
\label{sec:uncapped}

The uncapped rule has a separate obstruction: a large collection of
blocks shares one veto owner.  A fixed cap cannot simply be sent to
infinity in Theorem~\ref{thm:fixed-cap}, because doing so changes its
instance and its optimizing vector.

\begin{theorem}[Uncapped quadratic envelope]
\label{thm:uncapped}
For every $A>0$ and $N_0$ there is a binary finite-minlinear production
instance with $N\ge N_0$, $D=2N$, and a fixed full profile of output $1$
such that
\[
 E[q^{\infty,2}](I;\mathbf b)>A\sqrt N.
\]
Along an explicit unbounded sequence of owner counts,
\begin{equation}
 E[q^{\infty,2}](I;\mathbf b)>
 \frac1{40}\sqrt N\left(\frac{\log N}{\log\log N}\right)^{1/3}.
 \label{eq:uncapped-rate}
\end{equation}
The squared envelope norm and the ratio of envelope total to the full
seed total are also unbounded over these production instances.
\end{theorem}

\begin{lemma}[Full and focal base projections]
\label{lem:base-projections}
For the unit-demand row system of Lemma~\ref{lem:qary-input}, the full
quadratic projection is $z=k^{-1}\mathbf1$, with squared norm
$\rho=n/k^2$.  At the focal binary deletion, it is $w^i/t$, where
\begin{equation}
 t=1+(L-h)a+(m-K)b_*
  =1+(Q-1)L a^2+m b_*^2=\|w^i\|_2^2.
 \label{eq:base-t}
\end{equation}
Its focal coordinate, squared norm, and total row-multiplier mass are
all $\ell=1/t$, and $\ell\le\rho$.
\end{lemma}
\begin{proof}
Averaging any row over translations gives $(k/n)\mathbf1$.  Assign this
orbit distribution total multiplier $\rho$.  Its mean multiplier vector
is $z$, and every unit row is tight at $z$.  Thus for any feasible $y$,
$z\cdot y\ge\rho=\|z\|^2$, which proves that $z$ is the unique projection.

For the focal system use the mixture of Lemma~\ref{lem:qary-input}.
All supported rows have potential $t$ at $w^i$, and its displayed
coordinates give the second expression in \eqref{eq:base-t}.  Because
$a>b_*$, every off-focal row has potential at least
$La+(k-L)b_*$.  Its difference from $t$ is
\[
 h(a-b_*)-(1-b_*)>0.
\]
Consequently $w^i/t$ is feasible.  Give supported rows multipliers
$\alpha_R/t$, where $\alpha$ is the mixture.  Their sum is $1/t$ and
their weighted row vector is $w^i/t$.  The same projection inequality
proves uniqueness.  Deleting constraints only enlarges the feasible
set, so its squared minimum norm $\ell$ is at most $\rho$.
\end{proof}

Take $J$ disjoint copies of the base and adjoin the same owner $u$ to
every row.  At full demand the projection is
\begin{equation}
 x_u=\frac{J\rho}{1+J\rho},\qquad
 x^{(c)}=\frac{z}{1+J\rho}\quad(1\le c\le J).
 \label{eq:veto-full}
\end{equation}
Indeed, conditional on $x_u\in[0,1]$, each block is its full projection
scaled by $1-x_u$; minimize $x_u^2+(1-x_u)^2J\rho$.  Values $x_u>1$
cannot be optimal since $(1,0,\ldots,0)$ is feasible.

Delete the designated own rows in one focal block and leave all other
blocks unchanged.  With
$\eta=[1+(J-1)\rho+\ell]^{-1}$, the new projection is
\begin{equation}
 x_u=1-\eta,\qquad x^{(\rm focal)}=\eta w^i/t,
 \qquad x^{(c)}=\eta z\quad(c\ne\mathrm{focal}).
 \label{eq:veto-focal}
\end{equation}
To check this, scale the base row multipliers by $\eta$.  Their total
is $\eta((J-1)\rho+\ell)=1-\eta$, which is exactly the required veto
coordinate.  Retained supported rows are tight and all others are
feasible.  The resulting projection inequality proves optimality and
uniqueness.  In particular every private owner has selected own-report
payment
\begin{equation}
 p_J=\frac{\ell}{1+(J-1)\rho+\ell}.
 \label{eq:private-selected}
\end{equation}
Lemma~\ref{lem:realization} realizes all these reports with the same
$N=Jn+1$ actual owners and one production function on $2N$ resources.

\begin{proof}[Proof of Theorem~\ref{thm:uncapped}]
First take $j=1$ in Lemma~\ref{lem:qary-input} and put
$J=\lfloor1/\rho\rfloor$.  Every coordinate of $w^i$ is at least
$b_*>1/Q$, its coordinate sum is $k$, and
$t\le nb_*^2+QL$.  The exponential size of $n$ gives
$nb_*^2>QL$.  Hence
\[
 \rho<\frac{Q^2}{n}<\frac12,
 \qquad
 \frac\ell\rho=\frac{k^2}{nt}
       \ge\frac{nb_*^2}{nb_*^2+QL}>\frac12.
\]
Writing $c=J\rho$, we have $1/2<c\le1$ and $N\le2Jn$.  By
\eqref{eq:private-selected} and $\ell\le\rho$,
\begin{align*}
 \frac{Jn p_J}{\sqrt N}
 &\ge\frac{\ell\sqrt{Jn}}{2\sqrt2}
  =\frac{\ell}{\rho}\frac{\sqrt{c n\rho}}{2\sqrt2}
  >\frac{Q-1}{8},
\end{align*}
where $\sqrt{n\rho}=n/k>Q-1$.  Increasing $Q$ proves the first assertion.

For the rate, $L=16Q(Q-1)(Q+2)\le18Q^3$; the difference is
$2Q(Q-4)^2$.  Since $J\le k^2/n$, we have $n\le N\le k^2+1\le n^2$.
Therefore
\[
 \log N\le36Q^3\log Q,
 \qquad \log\log N>\log Q.
\]
Thus $(\log N/\log\log N)^{1/3}\le36^{1/3}Q$.
The preceding lower bound is greater than $Q\sqrt N/12$, and
$12\,36^{1/3}<40$, proving \eqref{eq:uncapped-rate}.

For the last two assertions, fix $Q$ and let $j\to\infty$ instead.
Put $b_0=1-\beta_0<1/(Q-1)$.  Then
$k/n\to b_0$, $t/n\to b_0^2$, $\ell/\rho\to1$,
$J\rho\to1$, and $N\sim b_0^2n^2$.  Equations
\eqref{eq:veto-full}--\eqref{eq:private-selected} imply
\[
 x_u^2+Jn p_J^2\longrightarrow\frac{1+b_0^{-2}}4,
 \qquad
 \frac{Jn p_J}{q^{\infty,2}(I;\mathbf b)}\longrightarrow\frac1{b_0}.
\]
These are lower bounds for the corresponding actual envelope quantities.
Letting $Q$ grow makes both unbounded.
\end{proof}

The first assertion and \eqref{eq:uncapped-rate} concern the true
all-real envelope, although their proofs only need selected admissible
underreports.  They also apply to every own-monotone rule preserving
the uncapped quadratic seed at every profile, by
Lemma~\ref{lem:least-majorant}.  No statement for every integer $N$ is
required or asserted.

\appendix
\section{The complete q-ary mixture construction}
\label{app:qary}

This appendix proves Lemma~\ref{lem:qary-input}.  No field structure is
used: the construction works for every integer $Q\ge3$.  Throughout,
color a vertex $x\in(\mathbb Z/Q\mathbb Z)^L$ by
$\sum_j x_j\pmod Q$, and write $M=Q^{L-1}$ for the size of each color
class.  The Hamming weight is the number of nonzero coordinates.
The focal vertex is $0$.  Its neighbors are all weight-one vertices.
The weight-two vertices form an exceptional set of size
\[
 e=(Q-1)^2\binom L2.
\]
Vertices of weight at least three will be called far vertices.  Let
$\mathcal F_c$ be the far vertices of color $c$, and $f_c=|\mathcal F_c|$.
Direct counting gives
\begin{equation}
 f_0=M-1-(Q-1)\binom L2,
 \qquad
 f_c=M-L-(Q-2)\binom L2\quad(c\ne0).
 \label{eq:far-counts}
\end{equation}
For a fixed pair of coordinates there are $Q-1$ nonzero coefficient
pairs summing to zero and $Q-2$ summing to any specified nonzero color.
This count also explains why composite $Q$ causes no exception.

\subsection{Uniform scalar inequalities}
\label{app:scalar}

The parameters of Lemma~\ref{lem:qary-input} satisfy $L\ge480$.
For every integer $D_0\ge20$,
\[
 3^{D_0-5}>100D_0^2.
\]
The assertion holds at $20$ and follows by induction because
$3D_0^2/(D_0+1)^2>1$.  Consequently
\begin{equation}
 M=Q^4 Q^{L-5}>100Q^4L^2.
 \label{eq:exponential-margin}
\end{equation}
In particular all far classes are nonempty, and
\begin{equation}
 M-QL^2\le f_c\le M,\qquad
 e+QL+2\le Q^2L^2,\qquad
 K-e\ge Q\beta_0 M-Q^2L^2,\qquad K\le Q\beta_0M.
 \label{eq:elementary-size-bounds}
\end{equation}
The last lower bound follows from $K\ge\beta_0m-1$ and
$m=QM-(Q-1)L-1$.

Define
\[
 B_{\rm av}=\frac{Q-1}{Q}-\frac{2\xi}{Q(Q-1)},\qquad
 \gamma_0=\frac{Q-1}{Q}-\beta_0.
\]
Substitution of $\xi=1/[2(Q+2)]$ gives
\begin{equation}
 B_{\rm av}=\beta_0+\delta,\qquad
 \gamma_0=\delta+\frac{2\xi}{Q(Q-1)}>\delta\ge\frac1{4Q^2}.
 \label{eq:availability-margin}
\end{equation}
Also $b_0=1-\beta_0<1/(Q-1)\le3/(2Q)$, so
$Q\beta_0-Q+2=2-Qb_0>1/2$.
The exponential bound implies
$m>h\ge2/\delta$, and hence
\begin{equation}
 \beta-\mu\ge\delta-1/m>\delta/2,\qquad
 h(\beta-\mu)-\beta
       \ge h\delta-h/m-\beta_0>0.
 \label{eq:off-margin-proof}
\end{equation}
Here $h\delta=2j\ge2$, $h/m<1$, and $\beta_0<1$.
Moreover $\beta\le\beta_0<(Q-1)/Q$.

For later transport, the following three inequalities hold for every
color $a_0$:
\begin{align}
 K-e&\ge (Q-2)M, \label{eq:hall-small}\\
 K&\le\sum_{c\ne a_0}f_c, \label{eq:hall-supply}\\
 K-e&\ge (Q\beta-Q+2)\sum_{c\ne a_0}f_c.
 \label{eq:hall-large}
\end{align}
For \eqref{eq:hall-small}, subtract $(Q-2)M$ in
\eqref{eq:elementary-size-bounds}; the difference is at least
$M/2-Q^2L^2>0$.  For \eqref{eq:hall-supply}, the difference between its
right and left sides is at least
$Q\gamma_0M-Q(Q-1)L^2>0$.  Finally $Q\beta-Q+2>0$, since
$\beta>\mu$ and
$Q\mu-Q+2=(Q-2+Q\xi)/(Q-1)>0$.
Replace $\beta$ by $\beta_0$ and the sum of far classes by $(Q-1)M$ to
bound the difference in \eqref{eq:hall-large} from below by
\[
 (Q-2)Q\gamma_0M-Q^2L^2>0.
\]
All strict inequalities follow from \eqref{eq:exponential-margin}.

\subsection{Local selection and repair}
\label{app:local-repair}

Choose an excluded color $c$ uniformly among the $Q$ colors.  Initially
retain all vertices of colors other than $c$; this is independent because
every complete axis line contains each color exactly once.

If $c=0$, keep a uniformly chosen subset of exactly $(Q-2)L+h$ of the
$(Q-1)L$ neighbors and leave all other initially retained vertices in
place.  If $c\ne0$, delete the focal vertex, choose a uniform $h$-subset
$S$ of coordinate indices, and add the previously excluded neighbors
$c e_j$ for $j\in S$.  For each added neighbor and each direction
$\ell\ne j$, independently choose $A$ uniformly from
$\{1,\ldots,Q-1\}$ and delete $c e_j+A e_\ell$.  Multiple deletions of
the same vertex have the effect of one deletion.

The resulting set is independent.  A line contains at most one added
neighbor.  The line through $c e_j$ in direction $j$ has lost the focal
vertex.  Every other line through it has lost the chosen repair vertex,
which was initially present because its color is $c+A\ne c$.
All remaining lines were independent before the additions.  The set
excludes $0$ and contains exactly $(Q-2)L+h$ neighbors; repairs only
delete weight-two vertices.

A neighbor $a_0e_j$ has conditional inclusion probability $\mu$ in phase
$c=0$, probability $\xi$ in phase $c=a_0$, and probability $1$ in the
other $Q-2$ phases.  Its unconditional marginal is therefore
$(\mu+\xi+Q-2)/Q=\mu$.

Consider a weight-two vertex $a_0e_j+b_0'e_\ell$, with nonzero
$a_0,b_0'$.  Its initial availability is $(Q-1)/Q$.  It can be deleted
from an added $a_0e_j$ in phase $a_0$, or from an added $b_0'e_\ell$ in
phase $b_0'$, each event having unconditional probability
$\xi/[Q(Q-1)]$.  Neither phase is its excluded color $a_0+b_0'$.
If the two phases agree, the union bound remains valid.  Thus its final
availability probability $P_v$ is rational and at least $B_{\rm av}$.
By \eqref{eq:availability-margin}, $P_v>\beta$.
Conditional on availability, retain this vertex with probability
$\beta/P_v$, independently for the different weight-two vertices.
Every exceptional vertex now has marginal exactly $\beta$.
This thinning preserves independence.  If $P$ is the number retained,
then $0\le P\le e$ and $\mathbb E P=\beta e$.

The remaining task is to select exactly $K-P$ far vertices with common
marginal $\beta$.  In phase $c$ all far vertices of colors other than
$c$ are still available, and none of color $c$ is available.  All local
choices and exceptional thinning have finite rational probabilities.
Let $\omega$ denote their complete state, of probability $p_\omega$,
with phase $c(\omega)$ and retained count $P_\omega$.  Set
\[
 A_\omega=\bigcup_{a_0\ne c(\omega)}\mathcal F_{a_0},
 \qquad T_\omega=K-P_\omega.
\]
By \eqref{eq:hall-small}--\eqref{eq:hall-supply},
$0\le T_\omega\le|A_\omega|$.  The total required marginal is correct:
\begin{equation}
 \mathbb E T_\omega=K-\beta e
                  =\beta(m-e)=\beta\Bigl|\bigcup_c\mathcal F_c\Bigr|.
 \label{eq:expected-supply}
\end{equation}

\subsection{Finite transport and all subset cuts}

\begin{lemma}[Finite fixed-cardinality transport]
\label{lem:transport}
Let $\mathcal F$ be finite.  In each of finitely many states $\omega$,
of rational probability $p_\omega>0$, let $A_\omega\subseteq\mathcal F$
be available and let $T_\omega$ be an integer between $0$ and
$|A_\omega|$.  Let $\beta$ be rational.  There is a finite rational
conditional distribution on $T_\omega$-subsets of $A_\omega$ giving
every vertex unconditional marginal $\beta$ if
\begin{equation}
 \mathbb E T_\omega=\beta|\mathcal F|,\qquad
 \beta|U|\le\sum_\omega p_\omega
             \min\{T_\omega,|U\cap A_\omega|\}\quad(U\subseteq\mathcal F).
 \label{eq:transport-cuts}
\end{equation}
\end{lemma}
\begin{proof}
Form a finite flow network with capacities $p_\omega T_\omega$ from
the source to each state, $p_\omega$ from that state to each available
vertex, and $\beta$ from each vertex to the sink.  For a fixed subset
$U$ of vertex nodes on the sink side, optimizing the placement of each
state gives cut capacity
\[
 \beta(|\mathcal F|-|U|)+
   \sum_\omega p_\omega\min\{T_\omega,|U\cap A_\omega|\}.
\]
Thus \eqref{eq:transport-cuts} and the finite max-flow/min-cut theorem
give a flow saturating every source and sink edge.  Capacities are
rational, so a rational flow can be chosen.  Dividing a state's outgoing
flows by $p_\omega$ gives a vector in $[0,1]^{A_\omega}$ with integer
sum $T_\omega$.

Such a rational vector is a rational convex combination of indicators
of $T_\omega$-subsets.  Indeed, if it is not integral it has at least
two fractional coordinates.  Perturb these coordinates in opposite
directions until one becomes integral at each endpoint.  The original
vector is a rational convex combination of the two endpoints; recursion
reduces the number of fractional coordinates.  This constructs the
required finite conditional distributions.
\end{proof}

We verify every cut in \eqref{eq:transport-cuts}.  For a subset $U$ of
far vertices, put $u_c=|U\cap\mathcal F_c|$.  The deficit
\[
 \sum_\omega p_\omega
   \min\left\{T_\omega,\sum_{c\ne c(\omega)}u_c\right\}
       -\beta\sum_c u_c
\]
is concave on the box $0\le u_c\le f_c$.  Every point of the box is a
convex combination of its corners, so it suffices to verify corners,
that is, unions of complete far color classes.

Suppose $U$ is the union of $r_0$ classes.  If $r_0\le Q-2$, then its
available count in every state is at most $(Q-2)M\le K-e\le T_\omega$.
The minimum is therefore the available count, whose expectation is
$(Q-1)|U|/Q\ge\beta|U|$.
If $r_0=Q-1$, let $a_0$ be the missing color.  In any phase belonging
to $U$, the available count is at most $(Q-2)M\le T_\omega$; averaging
these phases contributes $(Q-2)|U|/Q$.  In phase $a_0$, all of $U$ is
available and $T_\omega\le K\le|U|$, so the contribution is at least
$(K-e)/Q$.  Inequality \eqref{eq:hall-large} makes their sum at least
$\beta|U|$.  Finally, for all $Q$ classes the minimum is $T_\omega$
by \eqref{eq:hall-supply}, and equality follows from
\eqref{eq:expected-supply}.  The empty cut is immediate.

Lemma~\ref{lem:transport} now completes the far selection.  It only
removes vertices from the available independent set, so independence
is preserved.  Every realized set excludes $0$, has exactly
$(Q-2)L+h$ neighbors and exactly $K$ other vertices, and hence has size
$s$.  The neighbor and outside marginals are exactly $\mu$ and $\beta$.
Taking complements proves the mixture assertions in
Lemma~\ref{lem:qary-input}.  A row excluding $0$ is the complement of an
independent set containing $0$; each of the $L$ focal axis lines then
requires a neighbor in the row.  These neighbors are distinct.

It remains to check row size.  The lines parallel to one fixed axis
partition all vertices, and every independent set omits at least one
vertex from each; thus $k\ge n/Q$.  On the other hand, from the definition
of $K$,
\[
 k\le b_0n+\delta(Q-1)L+\beta_0+1
       < b_0n+\frac L{4Q}+2.
\]
Since $1/(Q-1)-b_0=\xi/(Q-1)+\delta>\delta$, the exponential estimate
gives $k<n/(Q-1)$.  The remaining inequalities in
\eqref{eq:qary-margins} are \eqref{eq:off-margin-proof} and
$\beta<(Q-1)/Q$.  Translating the construction supplies a mixture for
every focal vertex, completing the proof of
Lemma~\ref{lem:qary-input}.

\section{Simultaneous realization with fixed actual owners}
\label{app:realization}

The preceding row systems must be realized simultaneously.  Realizing
each underreport separately would not prove an envelope lower bound.
The next construction also shows why the production-resource count
differs from the local cube dimension.

\begin{lemma}[Binary incident deletions in one production function]
\label{lem:realization}
Let $I$ be a finite owner set with a distinguished owner $u$, and let
$\mathcal G$ be a nonempty finite family of subsets all containing $u$.
For each $j\in I$, prescribe a binary demand array
$d_G^j\in\{0,1\}$ satisfying $d_G^j=1$ whenever $j\notin G$;
take $d_G^u=1$ for every $G$.  Assume each array retains at least one
unit-demand generator.  There is a single binary finite-minlinear
production function on $2|I|$ resources and one fixed full bundle per
owner such that:
\begin{enumerate}
\item the full coalition game has the cover constraints
      $x(G)\ge1$ for $G\in\mathcal G$;
\item one designated own underreport by owner $j$ produces exactly the
      cover constraints $x(G)\ge d_G^j$;
\item the full and all designated underreport profiles have output $1$.
\end{enumerate}
The function is defined on all nonnegative real resource profiles, and
the owner set and full bundles are unchanged across the designated
underreports.
\end{lemma}
\begin{proof}
For $S\subseteq I$, define
\[
 v(S)=\max\bigl(\{1:G\in\mathcal G, G\subseteq S\}\cup\{0\}\bigr),
 \quad
 v^j(S)=\max\bigl(\{d_G^j:G\in\mathcal G, G\subseteq S\}\cup\{0\}\bigr).
\]
These monotone games take values in $[0,1]$, vanish without $u$, and
$v^j=v$ on coalitions not containing $j$.  For any such game $g$ and
nonempty $A\subseteq I$, define the nonnegative vector $h^{g,A}$ by
\[
 h_i^{g,A}=1\ (i\notin A),\qquad
 h_u^{g,A}=g(A)\ (u\in A),\qquad
 h_i^{g,A}=0\ (i\in A\setminus\{u\}).
\]
For every Boolean coalition vector,
\begin{equation}
 g(S)=\min_{\varnothing\ne A\subseteq I} h^{g,A}\cdot\mathbf1_S.
 \label{eq:game-forms}
\end{equation}
If $S\not\subseteq A$, the form contributes at least $1\ge g(S)$.
If $S\subseteq A$ and $u\in S$, it contributes $g(A)\ge g(S)$; if
$u\notin S$, the target is zero.  Taking $A=S$ gives equality for
nonempty $S$, and all forms vanish at the empty coalition.

Use resource coordinates $(X,Z)\in\mathbb R_+^I\times\mathbb R_+^I$ and
put
\begin{equation}
 F(X,Z)=\min\left\{
 h^{v,A}\cdot X,\quad h^{v^j,A}\cdot X+Z_j:
       \varnothing\ne A\subseteq I, j\in I\right\}.
 \label{eq:direct-production}
\end{equation}
All coefficients are binary.  Owner $j$ has the fixed full bundle
$(e_j,e_j)$.  At a full coalition, the baseline forms attain $v(S)$.
A branch $j$ with $j\in S$ has control contribution $1$ and cannot lower
that value; if $j\notin S$, its branch value is $v^j(S)=v(S)$.
Thus the full game is $v$.

If only owner $j$ hides its control resource, reporting $(e_j,0)$, its
branch attains $v^j(S)$.  The baseline and the branches of absent owners
have value $v(S)$, while all other present-owner branches have value at
least $1$.  The resulting game is exactly $v^j$.  Every one of these
games has grand-coalition value $1$ by the retained-generator assumption.
Finally, nonnegative payments cover a max-of-contained-generators game
if and only if they cover each generator at its stated demand: the
generator inequalities imply all coalition inequalities by containment,
and the converse follows by evaluating the game at a generator.
\end{proof}

For the q-ary constructions, the generators are $\{u\}\cup R$ in one
or more disjoint blocks, and the retained focal support is nonempty.
Lemma~\ref{lem:realization} therefore applies directly to every selected
binary deletion in the proofs above.  It establishes actual unilateral
reports, not changes to an abstract owner's identity or to the technology.
There is no assertion that the finite list of production forms is small
or efficiently enumerable.

There is also an efficient competing rule on the fixed full bundles'
downward box.  When $|I|\ge2$, the baseline form with $A=I\setminus\{u\}$
in \eqref{eq:direct-production} is $X_u$, so $F\le X_u$.  On that box,
only $u$ supplies this resource; every coalition without $u$ therefore
has value zero.  Paying the entire output to $u$ covers, is efficient,
and is own-monotone.  The additive-owner construction additionally pays
that owner's additive output.  These observations concern this box,
not arbitrary reallocations of resources among owners.  They reinforce
the distinction between failure of the specified baseline-preserving
repair and impossibility of an unrestricted mechanism.

\subsection{Every real own report preserves the generators}
\label{app:real-reports}

For the sharp class result below, selected reports alone are insufficient;
an upper bound must cover every real underreport.  We now verify that
\eqref{eq:direct-production} has exactly the needed property.  This
argument also works with prescribed arrays in $[0,1]$ instead of binary
arrays, although binary arrays suffice for our examples.

Fix a private owner $j\ne u$.  Every own underreport at its fixed full
bundle is $(t e_j,z e_j)$, with $t,z\in[0,1]$.  For a coalition $S$
containing $j$ and a game $g$ appearing in \eqref{eq:direct-production},
\begin{equation}
 \min_A h^{g,A}\cdot X_S
       =\min\{g(S),g(S\setminus\{j\})+t\}.
 \label{eq:real-base-form}
\end{equation}
If a member of $S\setminus\{j\}$ lies outside $A$, the form contributes
$1\ge g(S)$.  Otherwise it contains all those members.  Including $j$
minimizes at $A=S$, while excluding $j$ minimizes at
$A=S\setminus\{j\}$.  In the relevant positive-value case $u\in S$
and $u\ne j$, so this latter set is nonempty and is an allowed form.
If $u\notin S$, both sides have minimum zero, attained at $A=S$.
This handles singleton and empty boundary cases as well.

The other present owners' control branches cost at least $1$; an absent
owner's game agrees with $v$ on both sets in \eqref{eq:real-base-form}.
The exact new coalition value is consequently
\begin{equation}
 \min\{v(S),v(S\setminus\{j\})+t,v^j(S)+z\}
 \label{eq:real-coalition}
\end{equation}
when $j\in S$, and remains $v(S)$ otherwise.  To see that the extra
term $v^j(S\setminus\{j\})+t+z$ from its own branch is redundant, use
$v^j(S\setminus\{j\})=v(S\setminus\{j\})$ and $z\ge0$.
If $S$ contains an off-$j$ generator, \eqref{eq:real-coalition} is $1$.
If it contains no generator it is zero.  Otherwise all its generators
contain $j$, and finite maxima commute with the increasing truncation
in \eqref{eq:real-coalition}.  Thus the new game is exactly generated by
\begin{equation}
 d_G(t,z)=
 \begin{cases}
 1,&j\notin G,\\
 \min\{1,t,d_G^j+z\},&j\in G.
 \end{cases}
 \label{eq:real-generator-demands}
\end{equation}
Its cover polyhedron uses the same generators, with only incident demands
relaxed in $[0,1]$.

For a veto underreport $(t e_u,z e_u)$, the baseline form minimum on a
coalition is $t v(S)$.  A nonveto member outside $A$ costs $1$; exclusion
of $u$ costs $t$; and inclusion of all $S$ attains $t v(S)$.
Other present-owner branches cannot improve it, absent-owner branches
agree with $v$, and $v^u=v$.  Thus the whole game is exactly $tv$,
independently of $z$.  Its quadratic projection scales by $t$.

\section{The sharp fixed-arity common-veto class}
\label{app:sharp-class}

This appendix gives a matching upper and lower refinement for the
common-veto construction.  Its parameter $Q$ is the size of an
independence-forbidden hyperedge, not the production-resource dimension.
Fixing $Q$ restricts the construction class; the unrestricted inflation
in Theorem~\ref{thm:uncapped} lets $Q$ grow.

\subsection{A cap valid for all real own reports}

\begin{lemma}[Unit-demand multiplier mass]
\label{lem:dual-mass}
Let $A$ be a finite incidence matrix with nonempty rows, and suppose its
unit-demand cover is feasible.  Let $z$ be its minimum-norm cover and
$\rho=\|z\|^2$.  For any relaxed demands $d\in[0,1]^{\mathrm{rows}(A)}$,
the new projection $q$ has nonnegative row multipliers $\lambda$ with
$q=A^\mathsf T\lambda$.  Every such multiplier vector satisfies
\[
 q_i\le\sum_R\lambda_R\le\rho.
\]
At unit demand its total multiplier mass is $\rho$.
\end{lemma}
\begin{proof}
The optimality conditions for a quadratic objective on a polyhedron give
nonnegative row and nonnegativity multipliers.  At a positive coordinate
the latter vanish; at a zero coordinate nonnegativity of $A$ forces both
terms in the stationarity equation to vanish.  Thus $q=A^\mathsf T\lambda$.
Complementarity gives $\|q\|^2=\sum_R\lambda_R d_R\le L_*$, where
$L_*=\sum_R\lambda_R$.  For the full unit-demand problem these same row
multipliers have dual objective $2L_*-\|q\|^2\le\rho$.  Consequently
$L_*\le\rho$, and each coordinate of $q$ is a sum of a subset of the
multipliers.  At full demand complementarity gives $L_*=\rho$.
Alternatively, the dual inequality follows directly from
$\|y\|^2\ge2\lambda\cdot Ay-\|A^\mathsf T\lambda\|^2$ for any full
feasible $y$; no uniqueness of the multipliers is used.
\end{proof}

\begin{lemma}[Heterogeneous common-veto envelope bound]
\label{lem:veto-cap}
Take disjoint unit-demand base systems indexed by $c$, with $n_c$ owners,
full projections $z^c$, and $\rho_c=\|z^c\|^2$.  Add a common veto owner
to every row and let $R_* =\sum_c\rho_c$.  In the explicit production
realization \eqref{eq:direct-production}, at the fixed full bundles,
the true uncapped quadratic envelopes satisfy
\begin{equation}
 E_u=\frac{R_*}{1+R_*},\qquad
 E_i\le\frac{\rho_c}{1+R_*}\quad(i\text{ in block }c).
 \label{eq:class-cap}
\end{equation}
These inequalities also hold for the abstract supremum over all
incident generator-demand relaxations in $[0,1]$.
\end{lemma}
\begin{proof}
At full demand, the same minimization as in \eqref{eq:veto-full} gives
veto coordinate $R_* /(1+R_*)$ and squared norm $R_* /(1+R_*)$.
At any generator relaxation, if $\lambda$ are its row multipliers and
$u$ is its veto coordinate, stationarity gives $u=\sum\lambda$.
Complementarity and the full dual inequality give
$\|x\|^2\le u$ and $2u-\|x\|^2\le R_* /(1+R_*)$.
Hence $u\le R_* /(1+R_*)<1$.  The full demand attains this bound for the
veto.

For a private owner in block $c$, all other blocks remain unchanged.
Conditional on the new veto coordinate $u$, their projections are
$(1-u)z^d$.  Divide the focal block's coordinates and multipliers by
$1-u$.  This gives its base projection for residual demands
\[
 d'_G=\max\{0,(d_G-u)/(1-u)\}\in[0,1].
\]
Clipping nonpositive demands does not change the optimality conditions:
a nonempty row with positive multiplier has positive coordinate sum,
so cannot be tight at a nonpositive residual demand.  Let $L_*$ be the
normalized focal multiplier mass and $p_*$ its focal coordinate.
The veto stationarity equation is
$u=(1-u)(R_* -\rho_c+L_*)$.  Its actual private coordinate is therefore
\[
 \frac{p_*}{1+R_* -\rho_c+L_*}
 \le\frac{L_*}{1+R_* -\rho_c+L_*}
 \le\frac{\rho_c}{1+R_*},
\]
using Lemma~\ref{lem:dual-mass} and monotonicity in $L_*$.  Finally,
\eqref{eq:real-generator-demands} shows that every real private
underreport in the stated production realization is one of these
relaxations.  Every veto underreport scales the full game, as proved
in Appendix~\ref{app:real-reports}.  This proves the assertions for the
true envelope as well as the abstract relaxation envelope.
\end{proof}

\begin{theorem}[Sharp leading coefficient in the fixed-$Q$ class]
\label{thm:sharp-class}
Fix $Q\ge3$.  In each base block let the rows be a nonempty invariant
family of complements of fixed-cardinality independent sets in a
nonempty $Q$-uniform hypergraph.  Assume a vertex-transitive group
preserves both the hypergraph and the row family.  Form disjoint blocks
with a common veto and the direct production realization
\eqref{eq:direct-production}.  At its fixed full bundles, with output
$1$ and $N$ owners,
\begin{equation}
 E(I)\le\frac{1+\sqrt{1+Q^2(N-1)}}2,\qquad
 \sum_i E_i^2\le\frac{Q^4}{4(Q^2-1)}.
 \label{eq:class-finite-bounds}
\end{equation}
The supremum of attainable asymptotic leading funding coefficients is
$Q/2$.  More precisely, there are instances in this class with $N\to\infty$
and actual true envelopes satisfying
\begin{equation}
 \frac{E(I)}{\sqrt N}\longrightarrow\frac Q2,\qquad
 \sum_i E_i^2\longrightarrow\frac{Q^2+1}{4}.
 \label{eq:sharp-c-one}
\end{equation}
The squared-norm bound in \eqref{eq:class-finite-bounds} is itself
attained as a limit on a different sequence of block counts.
\end{theorem}
\begin{proof}[Upper bounds]
In a block of $n_c$ vertices every vertex has the same positive degree
$\Delta$.  If $W$ is an independent set, every edge contains a vertex
outside it.  Counting omissions gives
\[
 \Delta|V\setminus W|=\sum_{H\text{ an edge}}|H\setminus W|
                     \ge |\mathcal H|=n_c\Delta/Q.
\]
Thus the complement row size $k_c$ is at least $n_c/Q$.
Transitivity gives the uniform full projection and
$\rho_c=n_c/k_c^2$, exactly as in Lemma~\ref{lem:base-projections}.
Hence $n_c\rho_c\le Q^2$.

Put $m_* =\sum_c n_c=N-1$.  From \eqref{eq:class-cap},
\[
 E(I)\le\frac{R_*+\sum_c n_c\rho_c}{1+R_*}
 \le\frac{R_*+Q\sqrt{m_*R_*}}{1+R_*}.
\]
The second inequality uses $n_c\rho_c\le Q\sqrt{n_c\rho_c}$ and
Cauchy--Schwarz.  With $y=\sqrt{R_*}$ the last expression is the
Rayleigh quotient $(y^2+Q\sqrt{m_*}y)/(1+y^2)$, bounded by
$(1+\sqrt{1+Q^2m_*})/2$.  Similarly
\[
 \sum_i E_i^2\le
 \frac{R_*^2+\sum_c n_c\rho_c^2}{(1+R_*)^2}
 \le\frac{R_*^2+Q^2R_*}{(1+R_*)^2}
 \le\frac{Q^4}{4(Q^2-1)}.
\]
The last maximum occurs at $R_*=Q^2/(Q^2-2)$, as direct
differentiation shows.  This also proves the leading funding upper
coefficient $Q/2$.
\end{proof}

\begin{proof}[Construction attaining the limits]
We use the same local selection and finite transport as in
Appendix~\ref{app:qary}, with different parameters.  Fix $Q$ and let
\begin{align*}
 h&=8Q(Q-1),& L_0&=16Q(Q-1)(Q+3),& L&=jL_0\quad(j\ge1),\\
 n&=Q^L,& M&=Q^{L-1},&m&=n-(Q-1)L-1,\\
 \xi&=h/L,&\mu&=(Q-2+\xi)/(Q-1),&
 \beta_0&=(Q-1)/Q-24/L.
\end{align*}
Again put $K=\lfloor\beta_0m\rfloor$, $\beta=K/m$,
$s=(Q-2)L+h+K$, $k=n-s$, $a=1-\mu$, and $b_*=1-\beta$.
Take all complements of independent $s$-sets in the same axis-line
hypergraph.

We check every condition needed for the earlier distributional proof.
Its exceptional availability bound now is
\[
 B_{\rm av}=\frac{Q-1}{Q}-\frac{16}{L}
            =\beta_0+\frac8L>\beta,
\]
and
\[
 \beta_0-\mu=\frac1{Q(Q-1)}-\frac{8(Q+3)}L
             \ge\frac1{2Q(Q-1)}.
\]
The weight-two counts and far-class counts in
\eqref{eq:far-counts}--\eqref{eq:elementary-size-bounds} remain valid.
For every integer $D_0\ge20$, $3^{D_0-5}>100D_0^3$: check $D_0=20$
and use $3(D_0/(D_0+1))^3>1$.  Hence
\[
 M>100Q^4L^3.
\]
In particular $m>h$ and $\beta>\mu$.
Here $Q\beta_0-Q+2=1-24Q/L\ge7/8$ and
$Q-1-Q\beta_0=24Q/L$.  The three transport inequalities have the
following explicit positive margins:
\begin{align*}
 K-e-(Q-2)M&\ge(7/8)M-Q^2L^2>0,\\
 \sum_{c\ne a_0} f_c-K&\ge24QM/L-Q(Q-1)L^2>0,\\
 K-e-(Q\beta-Q+2)\sum_{c\ne a_0}f_c
       &\ge24Q(Q-2)M/L-Q^2L^2>0.
\end{align*}
The coefficient in the last line is positive since $\beta>\mu$.
Therefore all corner cuts, and by concavity all subset cuts, pass exactly
as in the proof of Lemma~\ref{lem:transport}.  Exceptional thinning,
the rational flow, and conditional fixed-cardinality selection produce
the same stated focal mixture for these parameters.

Its off-row potential gap is positive:
\[
 h(\beta-\mu)-\beta
       \ge4-h/m-\beta_0>2.
\]
Thus Lemma~\ref{lem:base-projections} applies with these new parameters,
giving the exact full norm $\rho=n/k^2$ and focal norm and multiplier
mass $\ell=1/t$.  As $j\to\infty$,
\[
 k/n\longrightarrow1/Q,\quad t/n\longrightarrow1/Q^2,\quad
 n\rho\longrightarrow Q^2,\quad \rho\longrightarrow0,\quad
 \ell/\rho\longrightarrow1.
\]
For any fixed $c_*>0$, choose integer block counts $J$ with
$J\rho\to c_*$.  The selected private payment
\eqref{eq:private-selected} and the true-envelope cap
\eqref{eq:class-cap} give, for every private owner,
\[
 \frac{\ell}{1+(J-1)\rho+\ell}
       \le E_i\le\frac{\rho}{1+J\rho}.
\]
The ratio of these two bounds tends to one, uniformly over owners.
The veto envelope is exactly $J\rho/(1+J\rho)$.
Since $N=Jn+1\sim c_*n^2/Q^2$, squeezing yields the actual limits
\[
 \frac{E(I)}{\sqrt N}\longrightarrow
         \frac{Q\sqrt{c_*}}{1+c_*},\qquad
 \sum_i E_i^2\longrightarrow
         \frac{c_*^2+Q^2c_*}{(1+c_*)^2}.
\]
Taking $c_*=1$ proves \eqref{eq:sharp-c-one}.  Taking
$c_*=Q^2/(Q^2-2)$ proves sharpness of the squared-norm bound.
The two optimizing scalings are different.  The same simultaneous
binary realization ensures that these limits concern the true
all-real envelopes for fixed actual owners in each instance.
\end{proof}

The nonempty-hypergraph assumption in Theorem~\ref{thm:sharp-class} is
essential for the positive-degree omission count.  Its upper bounds
also use the explicit generator-preserving realization; they are not
claims about every production function that happens to agree with a
given game at finitely many profiles.  Finally, fixing a single base
and merely sending its block count to infinity does not give the sharp
funding coefficient: the construction varies the local dimension while
choosing the block count so that $J\rho$ approaches the specified
positive constant.
